\documentclass[tikz,border=4pt]{standalone}
\usepackage{ctex}
\usepackage{amsmath}
\usetikzlibrary{calc, patterns}

\begin{document}
% 配方法几何意义：x^2 + bx + (b/2)^2 = (x + b/2)^2
% 用 b=6 即 x^2+6x+9=(x+3)^2 为例，几何拼接展示
\begin{tikzpicture}[scale=0.9, thick]

  % === 左侧：配方前 ===
  % x^2 正方形（边长 x，用 x=3 单位示意）
  \def\x{3}
  \def\b{1.5}  % b/2 = 3/2 = 1.5 单位

  \fill[blue!20] (0,0) rectangle (\x,\x);
  \draw (0,0) rectangle (\x,\x);
  \node at (\x/2, \x/2) {$x^2$};
  \node[below] at (\x/2, 0) {$x$};
  \node[left] at (0, \x/2) {$x$};

  % 上方 L 形：x * (b/2) 矩形
  \fill[red!20] (0,\x) rectangle (\x, \x+\b);
  \draw (0,\x) rectangle (\x, \x+\b);
  \node at (\x/2, \x+\b/2) {$\dfrac{b}{2}x$};
  \node[right] at (\x, \x+\b/2) {$\dfrac{b}{2}$};

  % 右方 L 形：(b/2) * x 矩形
  \fill[red!20] (\x,0) rectangle (\x+\b, \x);
  \draw (\x,0) rectangle (\x+\b, \x);
  \node at (\x+\b/2, \x/2) {$\dfrac{b}{2}x$};
  \node[above] at (\x+\b/2, \x) {$\dfrac{b}{2}$};

  % 缺角标注
  \draw[gray, dashed] (\x,\x) rectangle (\x+\b, \x+\b);
  \node[gray] at (\x+\b/2, \x+\b/2) {$\left(\dfrac{b}{2}\right)^{\!2}$};
  \node[gray, font=\small] at (\x+\b/2, \x+\b+0.4) {补入缺角};

  % 箭头和等式
  \draw[->, thick, blue] (\x+\b+0.5, \x/2+\b/2) -- (\x+\b+1.5, \x/2+\b/2)
      node[midway, above, font=\small] {$+\left(\dfrac{b}{2}\right)^2$};

  % === 右侧：配方后 ===
  \def\side{4.5}  % x+b/2 = 3+1.5 = 4.5
  \begin{scope}[xshift=7.2cm]
    \fill[blue!15] (0,0) rectangle (\side,\side);
    \draw[blue, line width=1.5pt] (0,0) rectangle (\side,\side);
    \node at (\side/2, \side/2) {$\left(x+\dfrac{b}{2}\right)^{\!2}$};
    \node[below] at (\side/2, 0) {$x + \dfrac{b}{2}$};
    \node[left] at (0, \side/2) {$x + \dfrac{b}{2}$};
  \end{scope}

  % 总说明
  \node[below] at (3.5, -0.6)
    {\small $x^2 + bx + \left(\dfrac{b}{2}\right)^2 = \left(x+\dfrac{b}{2}\right)^2$};

\end{tikzpicture}
\end{document}
